\subsection{A First Restriction on Matrix Outputs}

Suppose the second row is $\lambda$ times the first:
\[
A=\begin{pmatrix}a&b\\\lambda a&\lambda b\end{pmatrix}.
\]
For
\[
x=\begin{pmatrix}u\\v\end{pmatrix},
\]
we obtain
\[
Ax=
\begin{pmatrix}au+bv\\\lambda au+\lambda bv\end{pmatrix}
=
\begin{pmatrix}t\\\lambda t\end{pmatrix},
\qquad t=au+bv.
\]

\begin{interpretationbox}[The outputs satisfy a fixed relation]
No matter which input is chosen, the second output coordinate is $\lambda$ times the first. A zero determinant therefore signals that the possible outputs are restricted.
\end{interpretationbox}

\begin{examplebox}
For
\[
A=\begin{pmatrix}1&2\\2&4\end{pmatrix},
\]
$\det A=0$, and every output has the form
\[
Ax=\begin{pmatrix}t\\2t\end{pmatrix}.
\]
The outputs lie in one fixed pattern rather than filling the whole plane.
\end{examplebox}

A similar phenomenon occurs with three rows: if $r_3=r_1+r_2$, then the third entry of every output $Ax$ is the sum of the first two.
